\documentclass[12pt]{article}

% Idioma y codificación
\usepackage[spanish, es-tabla]{babel}       %es-tabla para que se titule "Tabla"
\usepackage[utf8]{inputenc}

% Márgenes
\usepackage[a4paper,top=3cm,bottom=2.5cm,left=3cm,right=3cm]{geometry}

% Comentarios de bloque
\usepackage{verbatim}

% Paquetes de links
\usepackage[hidelinks]{hyperref}    % Permite enlaces
\usepackage{url}                    % redirecciona a la web

% Más opciones para enumeraciones
\usepackage{enumitem}


%------------------------------------------------------------------------

%Paquetes de figuras
\usepackage{caption}
\usepackage{subcaption} % Figuras al lado de otras
\usepackage{float}      % Poner figuras en el sitio indicado H.


% Paquetes de imágenes
\usepackage{graphicx}       % Paquete para añadir imágenes
\usepackage{transparent}    % Para manejar la opacidad de las figuras

% Paquete para usar colores
\usepackage[dvipsnames]{xcolor}
\usepackage{pagecolor}      % Para cambiar el color de la página

% Habilita tamaños de fuente mayores
\usepackage{fix-cm}


%------------------------------------------------------------------------

% Paquetes de matemáticas
\usepackage{amsmath, amsfonts, amssymb, mathrsfs}
\usepackage{cancel}     % Simplificar tachando
\usepackage{polynom}    % Divisiones y Ruffini

\usepackage{pgfplots}   %Representar funciones
\pgfplotsset{compat=1.18}  % Versión 1.18

%Definición de teoremas, etc.
\usepackage{amsthm}
%\swapnumbers   % Intercambia la posición del texto y de la numeración

\theoremstyle{plain}
\newtheorem{teo}{Teorema}[chapter]     % Se resetea en cada chapter
\newtheorem{coro}{Corolario}[teo]           % Se resetea en cada teorema
\newtheorem{prop}[teo]{Proposición}         % Usa el mismo contador que teorema
\newtheorem{lema}[teo]{Lema}                % Usa el mismo contador que teorema

\theoremstyle{remark}
\newtheorem*{observacion}{Observación}

\theoremstyle{definition}
\newtheorem{definicion}{Definición} [chapter]     % Se resetea en cada chapter
\newtheorem*{notacion}{Notacion}
\newtheorem*{ejemplo}{Ejemplo}
\newtheorem*{ejercicio*}{Ejercicio}             % No numerado
\newtheorem{ejercicio}{Ejercicio} [section]     % Se resetea en cada section


% Modificar el formato de la numeración del teorema "ejercicio"
\renewcommand{\theejercicio}{%
  \ifnum\value{section}=0 % Si no se ha iniciado ninguna sección
    \arabic{ejercicio}% Solo mostrar el número de ejercicio
  \else
    \thesection.\arabic{ejercicio}% Mostrar número de sección y número de ejercicio
  \fi
}


% \renewcommand\qedsymbol{$\blacksquare$}         % Cambiar símbolo QED
%------------------------------------------------------------------------

% Paquetes para encabezados
\usepackage{fancyhdr}
\pagestyle{fancy}
\fancyhf{}

\newcommand{\helv}{ % Modificación tamaño de letra
\fontfamily{}\fontsize{12}{12}\selectfont}
\setlength{\headheight}{15pt} % Amplía el tamaño del índice


%\usepackage{lastpage}   % Referenciar última pag   \pageref{LastPage}
\fancyfoot[C]{\thepage}

%------------------------------------------------------------------------

\begin{comment} % Descomentar en book. Comentar en article

% Conseguir que no ponga "Capítulo 1". Sino solo "1."
\makeatletter
\renewcommand{\@makechapterhead}[1]{%
  \vspace*{50\p@}%
  {\parindent \z@ \raggedright \normalfont
    \ifnum \c@secnumdepth >\m@ne
      \huge\bfseries \thechapter.\hspace{1em}\ignorespaces
    \fi
    \interlinepenalty\@M
    \Huge \bfseries #1\par\nobreak
    \vskip 40\p@
  }}
\makeatother
\end{comment}

%------------------------------------------------------------------------
% Personalizar la portada
\usepackage{titling}


%------------------------------------------------------------------------

% Comandos definidos
\newcommand{\bb}[1]{\mathbb{#1}}
\newcommand{\cc}[1]{\mathcal{#1}}
\newcommand{\Hop}{\stackrel{L'H\hat{o}pital}{=}}

\newcommand{\red}[1]{{\color{red}#1}} % Para integrales, destacar los cambios.

% Método de integración
\newcommand{\MetInt}[2]{
    \left[\begin{array}{c}
        #1 \\ #2
    \end{array}\right]
}

%------------------------------------------------------------------------
\usepackage{extarrows}
\usepackage{stackrel}
\usetikzlibrary{matrix} % Para divisiones de polinomios.

% En el preámbulo del documento o en tu archivo .sty
\newcounter{ejercicio}[section] % Define el contador de ejercicio y lo reinicia con cada sección

\newcounter{ejercicio}
\newcommand{\resetearcontador}{%
  \setcounter{ejercicio}{0} % Resetea el contador de ejercicios a 0
}

\renewcommand{\theejercicio}{\arabic{ejercicio}}

\begin{document}

    % 1. Foto de fondo
    % 2. Título
    % 3. Encabezado Izquierdo
    % 4. Color de fondo
    % 5. Coord x del titulo
    % 6. Coord y del titulo
    % 7. Fecha

    
    % 1. Foto de fondo
% 2. Título
% 3. Encabezado Izquierdo
% 4. Color de fondo
% 5. Coord x del titulo
% 6. Coord y del titulo
% 7. Fecha

\newcommand{\portada}[7]{

    \portadaBase{#1}{#2}{#3}{#4}{#5}{#6}{#7}
    \portadaBook{#1}{#2}{#3}{#4}{#5}{#6}{#7}
}

\newcommand{\portadaExamen}[7]{

    \portadaBase{#1}{#2}{#3}{#4}{#5}{#6}{#7}
    \portadaArticle{#1}{#2}{#3}{#4}{#5}{#6}{#7}
}




\newcommand{\portadaBase}[7]{

    % Tiene la portada principal y la licencia Creative Commons
    
    % 1. Foto de fondo
    % 2. Título
    % 3. Encabezado Izquierdo
    % 4. Color de fondo
    % 5. Coord x del titulo
    % 6. Coord y del titulo
    % 7. Fecha
    
    
    \thispagestyle{empty}               % Sin encabezado ni pie de página
    \newgeometry{margin=0cm}        % Márgenes nulos para la primera página
    
    
    % Encabezado
    \fancyhead[L]{\helv #3}
    \fancyhead[R]{\helv \nouppercase{\leftmark}}
    
    
    \pagecolor{#4}        % Color de fondo para la portada
    
    \begin{figure}[p]
        \centering
        \transparent{0.3}           % Opacidad del 30% para la imagen
        
        \includegraphics[width=\paperwidth, keepaspectratio]{assets/#1}
    
        \begin{tikzpicture}[remember picture, overlay]
            \node[anchor=north west, text=white, opacity=1, font=\fontsize{60}{90}\selectfont\bfseries\sffamily, align=left] at (#5, #6) {#2};
            
            \node[anchor=south east, text=white, opacity=1, font=\fontsize{12}{18}\selectfont\sffamily, align=right] at (9.7, 3) {\textbf{\href{https://losdeldgiim.github.io/}{Los Del DGIIM}}};
            
            \node[anchor=south east, text=white, opacity=1, font=\fontsize{12}{15}\selectfont\sffamily, align=right] at (9.7, 1.8) {Doble Grado en Ingeniería Informática y Matemáticas\\Universidad de Granada};
        \end{tikzpicture}
    \end{figure}
    
    
    \restoregeometry        % Restaurar márgenes normales para las páginas subsiguientes
    \pagecolor{white}       % Restaurar el color de página
    
    
    \newpage
    \thispagestyle{empty}               % Sin encabezado ni pie de página
    \begin{tikzpicture}[remember picture, overlay]
        \node[anchor=south west, inner sep=3cm] at (current page.south west) {
            \begin{minipage}{0.5\paperwidth}
                \href{https://creativecommons.org/licenses/by-nc-nd/4.0/}{
                    \includegraphics[height=2cm]{assets/Licencia.png}
                }\vspace{1cm}\\
                Esta obra está bajo una
                \href{https://creativecommons.org/licenses/by-nc-nd/4.0/}{
                    Licencia Creative Commons Atribución-NoComercial-SinDerivadas 4.0 Internacional (CC BY-NC-ND 4.0).
                }\\
    
                Eres libre de compartir y redistribuir el contenido de esta obra en cualquier medio o formato, siempre y cuando des el crédito adecuado a los autores originales y no persigas fines comerciales. 
            \end{minipage}
        };
    \end{tikzpicture}
    
    
    
    % 1. Foto de fondo
    % 2. Título
    % 3. Encabezado Izquierdo
    % 4. Color de fondo
    % 5. Coord x del titulo
    % 6. Coord y del titulo
    % 7. Fecha


}


\newcommand{\portadaBook}[7]{

    % 1. Foto de fondo
    % 2. Título
    % 3. Encabezado Izquierdo
    % 4. Color de fondo
    % 5. Coord x del titulo
    % 6. Coord y del titulo
    % 7. Fecha

    % Personaliza el formato del título
    \pretitle{\begin{center}\bfseries\fontsize{42}{56}\selectfont}
    \posttitle{\par\end{center}\vspace{2em}}
    
    % Personaliza el formato del autor
    \preauthor{\begin{center}\Large}
    \postauthor{\par\end{center}\vfill}
    
    % Personaliza el formato de la fecha
    \predate{\begin{center}\huge}
    \postdate{\par\end{center}\vspace{2em}}
    
    \title{#2}
    \author{\href{https://losdeldgiim.github.io/}{Los Del DGIIM}}
    \date{Granada, #7}
    \maketitle
    
    \tableofcontents
}




\newcommand{\portadaArticle}[7]{

    % 1. Foto de fondo
    % 2. Título
    % 3. Encabezado Izquierdo
    % 4. Color de fondo
    % 5. Coord x del titulo
    % 6. Coord y del titulo
    % 7. Fecha

    % Personaliza el formato del título
    \pretitle{\begin{center}\bfseries\fontsize{42}{56}\selectfont}
    \posttitle{\par\end{center}\vspace{2em}}
    
    % Personaliza el formato del autor
    \preauthor{\begin{center}\Large}
    \postauthor{\par\end{center}\vspace{3em}}
    
    % Personaliza el formato de la fecha
    \predate{\begin{center}\huge}
    \postdate{\par\end{center}\vspace{5em}}
    
    \title{#2}
    \author{\href{https://losdeldgiim.github.io/}{Los Del DGIIM}}
    \date{Granada, #7}
    \thispagestyle{empty}               % Sin encabezado ni pie de página
    \maketitle
    \vfill
}
    \portadaExamen{ffccA4.jpg}{Álgebra I\\Examen III}{Álgebra I. Examen III}{MidnightBlue}{-8}{28}{2023-2024}{José Juan Urrutia Milán\\Arturo Olivares Martos}

    
    \begin{description}
        \item[Asignatura] Álgebra I.
        \item[Curso Académico] 2021-22.
        \item[Grado] Doble Grado en Ingeniería Informática y Matemáticas.
        \item[Grupo] Único.
        \item[Profesor] María Pilar Carrasco Carrasco.
        \item[Descripción] Examen Extraordinario.
        \item[Fecha] 7 de febrero de 2022.
        \item[Duración] 3 horas.
    
    \end{description}
    \newpage
    
    \begin{ejercicio}[4 puntos]
        \ 
        \begin{enumerate} 
            \item (0.5 puntos) Sea $X$ un conjunto y $A$, $B$ y $C$ subconjuntos de $X$. Suponiendo que $A \cap C = \emptyset$, probad que $(A \cup B)-C=A\cup (B-C)$.
            \item (1 punto) Sea $f:\bb{R}\rightarrow \bb{R}$ la aplicación definida para cada $x \in \bb{R}$, por $f(x)~=~x^2-3x+1$. Sea $R_f$ la relación de equivalencia en $\bb{R}$ definida por la aplicación $f$. Describid el conjunto cociente. ¿Cuántos elementos tiene cada clase?
            \item (0.5 puntos) En el anillo $\bb{Z}_8$, sea $I = 4\bb{Z}_8$, el ideal principal generado por 4. Describid $\bb{Z}_8/I$ listando todos sus elementos.
            \item (1 punto) Sea $A$ un dominio euclídeo con función euclídea $\phi:A \setminus \{0\}\rightarrow \bb{N}$. Demostrad que para todo $a \neq 0$, se tiene que $\phi(a) \geq \phi(1)$ y que se da la igualdad si, y sólo si $a \in U(A)$.
            \item (0.5 puntos) Demostrad que $n^{13}-n$ es divisible por 2 y 5 para todo $n \in \bb{Z}$.
            \item (0.5 puntos) Demostrad que en $\bb{Z}[i]$ se tiene que mcd$(n,n+i) = 1$, para todo $n \in \bb{Z}$, $n \neq 0$.
        \end{enumerate} 
    \end{ejercicio}
    
    \begin{ejercicio}[3.5 puntos]
        \ 
        \begin{enumerate}
            \item (2.5 puntos) Factorizar en $\bb{Z}[x]$ y en $\bb{Q}[x]$ los siguientes polinomios:
            \begin{enumerate}
                \item[a.] $48x^4 + 24x^3 -72x +80$
                \item[b.] $x^6 + 8x^4 + 4x^3 + 5x^2 + 4x+1$
                \item[c.] $10x^5 + 23x^4 -20x^3 + 5x^2 + 5x-3$
            \end{enumerate}
            \item Sea $f(x) \in \bb{Z}[x]$ un polinomio primitivo de grado $n>0$ y tal que existe un primo $p \in \bb{Z}$ verificando:
                \begin{enumerate}
                    \item[(i)] Su reducido módulo $p$ es de la forma $R_p(f(x)) = \alpha x^n$, con $\alpha \in \bb{Z}$ no nulo.
                    \item[(ii)] $p^2 \nmid f(p)$. 
                \end{enumerate}
                Demostrad que $f(x)$ es irreducible.
        \end{enumerate}
    \end{ejercicio}

    \begin{ejercicio}[2.5 puntos]
        Dado el sistema de congruencias en $\bb{Z}[\sqrt{-2}]$
        $$\left\{ \begin{tabular}{l}
            x \equiv 1 \mod (1+\sqrt{-2}) \\
            x \equiv 2 \mod (3+\sqrt{-2}) \\
            x \equiv 4 \mod (3+2\sqrt{-2}) 
        \end{tabular}\right.$$
        Demostrad que tiene solución sin resolverlo. Resolverlo, dando una solución óptima y la solución general. \newline
        ¿Es posible encontrar una solución a $a + b\sqrt{-2}$ tal que $30 < a < 56$?
    \end{ejercicio}

    \newpage
    \ % --------------------------------------------------------------------------------
    \newpage
    \resetearcontador
    
    \begin{ejercicio}[4 puntos]
        \ 
        \begin{enumerate} 
            \item (0.5 puntos) Sea $X$ un conjunto y $A$, $B$ y $C$ subconjuntos de $X$. Suponiendo que $A \cap C = \emptyset$, probad que $(A \cup B)-C=A\cup (B-C)$.
            \item (1 punto) Sea $f:\bb{R}\rightarrow \bb{R}$ la aplicación definida para cada $x \in \bb{R}$, por $f(x)~=~x^2-3x+1$. Sea $R_f$ la relación de equivalencia en $\bb{R}$ definida por la aplicación $f$. Describid el conjunto cociente. ¿Cuántos elementos tiene cada clase?
            \item (0.5 puntos) En el anillo $\bb{Z}_8$, sea $I = 4\bb{Z}_8$, el ideal principal generado por 4. Describid $\bb{Z}_8/I$ listando todos sus elementos.\\

                \noindent
                $$\bb{Z}_8/I = \{[a] \mid a \in \bb{Z}_8\} = \{a+I \mid a \in \bb{Z}_8\}$$
                $$[0] = \{ b \in \bb{Z}_8 \mid b \equiv 0 \mod (4) \} = \{ b \in \bb{Z}_8 \mid b = 0 \mbox{ en } \bb{Z}_4 \} = \{ 0,4 \}$$
                $$[1] = \{ b \in \bb{Z}_8 \mid b \equiv 1 \mod (4) \} = \{ b \in \bb{Z}_8 \mid b = 1 \mbox{ en } \bb{Z}_4 \} = \{ 1,5 \}$$
                $$[2] = \{ b \in \bb{Z}_8 \mid b \equiv 2 \mod (4) \} = \{ b \in \bb{Z}_8 \mid b = 2 \mbox{ en } \bb{Z}_4 \} = \{ 2,6 \}$$
                $$[3] = \{ b \in \bb{Z}_8 \mid b \equiv 3 \mod (4) \} = \{ b \in \bb{Z}_8 \mid b = 3 \mbox{ en } \bb{Z}_4 \} = \{ 3,7 \}$$
                Luego:
                $$\bb{Z}_8/I = \{ 0+I, 1+I, 2+I, 3+I \}$$

            \item (1 punto) Sea $A$ un dominio euclídeo con función euclídea $\phi:A \setminus \{0\}\rightarrow \bb{N}$. Demostrad que para todo $a \neq 0$, se tiene que $\phi(a) \geq \phi(1)$ y que se da la igualdad si, y sólo si $a \in U(A)$.\\

                $$\phi(a) = \phi(a \cdot 1) \geq \phi(1)~~\forall a \in A \setminus \{0\}$$

                \noindent
                $\Rightarrow)$ Supongamos que $\phi(a) = \phi(1)$:\newline
                Dividimos 1 entre $a$, $1 = qa + r$ con:
                $$\left\{ \begin{tabular}{l}
                    r = 0 \\
                    \lor \\
                    \phi(r) < \phi(a) = \phi(1)
                \end{tabular}\right.$$
                $\bullet$ Supongamos que $r \neq 0$, luego $\phi(r) < \phi(1)$, pero:
                $$\phi(b) \geq \phi(1)~~\forall b \in A \setminus\{0\} \Rightarrow \phi(r) \geq \phi(1)$$
                \underline{Contradicción}. Luego $r = 0$.\newline
                Por tanto: $1 = qa \Rightarrow a \in U(A)$.\\

                \noindent
                $\Leftarrow)$ Supongamos que $a \in U(A) \Rightarrow \exists a^{-1} \in A \mid aa^{-1} = 1 \Rightarrow \phi(aa^{-1}) = \phi(1)$
                $$\left. \begin{tabular}{r}
                    \phi(1) = \phi(aa^{-1}) \geq \phi(a) \\
                    \phi(a) \geq \phi(1)~~\forall a \in A \setminus\{0\}
                \end{tabular}\right\}\Rightarrow \phi(a) = \phi(1)$$

            \item (0.5 puntos) Demostrad que $n^{13}-n$ es divisible por 2 y 5 para todo $n \in \bb{Z}$.\\

                \noindent
                Demostremos primero que $2 \mid n^{13}-n~~\forall n \in \bb{Z}$:\newline
                $\bullet$ Supongamos que $2 \mid n \Rightarrow \exists k \in \bb{Z}$ tal que $n = 2k$.
                $$n^{13}-n = n(n^{12}-1) = 2k(n^{12}-1)=2k'$$
                Con $k' = k(n^{12}-1) \Rightarrow 2 \mid n^{13}-n$.

                \noindent
                $\bullet$ Supongamos que $2 \nmid n$:\newline
                Como 2 es primo, mcd$(2,n) = 1$. Por el Teorema de Fermat:
                $$n^{\varphi(2)} = n \equiv 1 \mod (2) \Rightarrow n^{12} \equiv 1 \mod (2) \Rightarrow n^{13} \equiv n \mod (2) \Rightarrow$$
                $$\Rightarrow n^{13}-n \equiv 0 \mod (2) \Rightarrow 2 \mid n^{13}-n$\\

                \noindent
                Veamos que $5 \mid n^{13}-n~~\forall n \in \bb{Z}$:\newline
                $\bullet$ Supongamos que $5 \mid n \Rightarrow \exists k \in \bb{Z}$ tal que $n = 5k$.
                $$n^{13}-n = n(n^{12}-1) = 5k(n^{12}-1)=5k'$$
                Con $k' = k(n^{12}-1) \Rightarrow 5 \mid n^{13}-n$.

                \noindent
                $\bullet$ Supongamos que $5 \nmid n$:\newline
                Como 5 es primo, mcd$(5,n) = 1$. Por el Teorema de Fermat:
                $$n^{\varphi(5)} = n^4 \equiv 1 \mod (5) \Rightarrow n^{12} \equiv 1 \mod (5) \Rightarrow n^{13} \equiv n \mod (5) \Rightarrow$$
                $$\Rightarrow n^{13}-n \equiv 0 \mod (5) \Rightarrow 5 \mid n^{13}-n$$


            \item (0.5 puntos) Demostrad que en $\bb{Z}[i]$ se tiene que mcd$(n,n+i) = 1$, para todo $n \in \bb{Z}$, $n \neq 0$.\\

                \noindent
                Supongamos que $\exists n \in \bb{Z}^{+} \mid$ mcd$(n, n+i) \neq 1$.\newline
                Como el máximo común divisor es único salvo asociados y:\newline $1\sim a$ con $a \in U(\bb{Z}[i])$:
                $$\mbox{mcd}(n, n+i) = \alpha \mid \alpha \notin U(A)$$
                $$\left\{ \begin{tabular}{lclcl}
                    \alpha \mid n & \Rightarrow & N(\alpha) \mid N(n) & \Rightarrow & N(\alpha) \mid n^2 \\
                    \land & & & & \\
                    \alpha \mid n+i & \Rightarrow & N(\alpha) \mid N(n+i) & \Rightarrow & N(\alpha) \mid n^2 +1
                \end{tabular}\right.$$
                
                Veamos ahora el siguiente resultado:\newline
                Sean $a$, $b \in \bb{Z}^{+}$: $a \mid b \land a\mid b+1 \Leftrightarrow a = \pm 1$

                \noindent
                $\Leftarrow)$ Supongamos que $a = \pm 1 \Rightarrow a \in U(A) \Rightarrow a \mid b \land a \mid b+1$\newline
                $\Rightarrow)$ Supongamos que $a \mid b \land a \mid b + 1$:
                $$\left. \begin{tabular}{lcr}
                    a \mid b & \Rightarrow & \exists k \in \bb{Z} \mid b = ka \\
                    a \mid b +1 & \Rightarrow & \exists k' \in \bb{Z} \mid b+1 = k'a
                \end{tabular}\right\} \Rightarrow ka+1 = k'a \Rightarrow$$
                $$\Rightarrow 1 = a(k'-k) \Rightarrow a \mid 1 \Rightarrow a \in U(\bb{Z}) \Rightarrow a = \pm 1$$

                $$\left. \begin{tabular}{l}
                    N(\alpha) \mid n^2 \\
                    \land \\
                    N(\alpha) \mid n^2 + 1
                \end{tabular}\right\} \Rightarrow N(\alpha) = \pm 1 \Rightarrow \alpha \in U(A)$$
                Lo que es una \underline{contradicción}.

                \noindent
                Luego $mcd(n, n+i) = 1~~\forall n \in \bb{Z}^{+}$.
        \end{enumerate} 
    \end{ejercicio}
    
    \begin{ejercicio}[3.5 puntos]
        \ 
        \begin{enumerate}
            \item (2.5 puntos) Factorizar en $\bb{Z}[x]$ y en $\bb{Q}[x]$ los siguientes polinomios:
            \begin{enumerate}
                \item[a.] $48x^4 + 24x^3 -72x +80$\\

                    \noindent
                    Sea $f =48x^4 + 24x^3 -72x +80 $, $f=8g$ con $g = 6x^4+3x^3-9x+10 \in \bb{Z}[x]$.\newline
                    $g$ es irreducible en $\bb{Z}[x]$ por Eisenstein para $p = 3$. Por la misma razón, también lo es en $\bb{Q}[x]$.
                    $$f = 2^3 (6x^4+3x^3-9x+10 ) \mbox{ en } \bb{Z}[x]$$
                    $$f = 8 (6x^4+3x^3-9x+10 ) \mbox{ en } \bb{Q}[x]$$

                \item[b.] $x^6 + 8x^4 + 4x^3 + 5x^2 + 4x+1$\\

                    \noindent
                    Sea $f =x^6 + 8x^4 + 4x^3 + 5x^2 + 4x+1 $.\newline
                    Las posibles raíces de $f$ en $\bb{Q}$ son $\{\pm 1\}$:
                    $$\left.\begin{tabular}{l}
                            $f$(1) = 1 + 8 + 4 + 5 + 4 + 1 = 23 \neq 0 \\
                            $f$(-1) = 1 + 8 -4 + 5 -4 + 1 = 7 \neq 0
                        \end{tabular}\right\} \Rightarrow \begin{tabular}{l}
                        $f$ \mbox{ no tiene factores } \\
                        \mbox{ de grado 1 ni 5 }
                    \end{tabular}$$
                    $\bullet$ Reducimos módulo 2:
                    $$R_2(f) = x^6 +x^2 + 1 = (x^3 +x +1 )^3 \Rightarrow R_2(f) \mbox{ no tiene factores de grado } 2 \mbox{ ni } 4$$
                    Luego $f$ tampoco tiene factores de grado 2 ni 4.\\

                    \noindent
                    $\bullet$ Reducimos módulo 3:
                    $$R_3(f) = x^6 + 2x^4 + x^3 + 2x^2 + x + 1 \in \bb{Z}_3[x]$$
                    $$\left. \begin{tabular}{l}
                        $R_3(f)(0) = 1\neq 0$ \\
                        $R_3(f)(1) = 8 \neq 0$ \\
                        $R_3(f)(2) = 115 = 1 \neq 0$ \\
                    \end{tabular}\right\} \Rightarrow \begin{tabular}{l}
                        $R_3(f)$ \mbox{ no tiene factores } \\
                        \mbox{ de grado 1 ni 5 }
                    \end{tabular}$$
                    
                    Al dividir $R_3(f)$ entre $x^2+1$ obtenemos que:
                    $$x^6 + 2x^4 + x^3 + 2x^2 + x + 1 = (x^2+1)(x^4+x^2+x+1) \Rightarrow (x^2+1) \mid R_3(f)$$
                    $$R_3(f) = (x^2+1)g \mbox{ con } g=x^4 +x^2+x+1 \in \bb{Z}_3[x]$$
                    Como $R_3(f)$ no tiene factores de grado 1 $\Rightarrow g$ no tiene de grado 1 ni 3.\newline
                    $R_3(f)$ no tiene factores de grado 3 $\Rightarrow f$ tampoco.\\

                    \noindent
                    Concluimos que $f$ es irreducible en $\bb{Z}[x]$ y en $\bb{Q}[x]$.


                \item[c.] $10x^5 + 23x^4 -20x^3 + 5x^2 + 5x-3$\\

                    \noindent
                    

            \end{enumerate}
            \item Sea $f(x) \in \bb{Z}[x]$ un polinomio primitivo de grado $n>0$ y tal que existe un primo $p \in \bb{Z}$ verificando:
                \begin{enumerate}
                    \item[(i)] Su reducido módulo $p$ es de la forma $R_p(f(x)) = \alpha x^n$, con $\alpha \in \bb{Z}$ no nulo.
                    \item[(ii)] $p^2 \nmid f(p)$. 
                \end{enumerate}
                Demostrad que $f(x)$ es irreducible.
        \end{enumerate}
    \end{ejercicio}

    \begin{ejercicio}[2.5 puntos]
        Dado el sistema de congruencias en $\bb{Z}[\sqrt{-2}]$
        $$\left\{ \begin{tabular}{l}
            x \equiv 1 \mod (1+\sqrt{-2}) \\
            x \equiv 2 \mod (3+\sqrt{-2}) \\
            x \equiv 4 \mod (3+2\sqrt{-2}) 
        \end{tabular}\right.$$
        Demostrad que tiene solución sin resolverlo. Resolverlo, dando una solución óptima y la solución general. \newline
        ¿Es posible encontrar una solución a $a + b\sqrt{-2}$ tal que $30 < a < 56$?
    \end{ejercicio}


\end{document}
